% ============================================================
% section3_loading.tex - Раздел 3: Загрузка и первичная визуализация данных
% ============================================================
% !TeX program = xelatex
% Компиляция: xelatex --shell-escape main.tex

\section{Загрузка и первичная визуализация экспериментальных данных}

Прежде чем обрабатывать сигналы, нужно научиться правильно их читать и отображать. Экспериментальные данные — это фундамент, на котором строится вся дальнейшая физика. В этом разделе мы разберемся, как прибор записывает результаты, научимся извлекать из имен файлов параметры измерений, сразу же построим наши первые графики с помощью библиотеки Matplotlib и спроектируем правильную процедуру загрузки всех файлов в единое хранилище (базу данных).

\subsection{Формат экспериментальных данных}

Наши экспериментальные данные получены на установке импульсной терагерцовой спектроскопии (Time-Domain Spectroscopy, TDS). В такой системе измеряется временной профиль напряженности электрического поля ТГц-импульса $E(t)$.

Вот в каком формате записаны текстовые файлы измерений в папке \texttt{data}:
\begin{itemize}
    \item Внутри каждого файла — \textbf{три колонки чисел}, разделенные табуляцией: время задержки (в пикосекундах), амплитуда ТГц-сигнала (в условных единицах) и текущая позиция линии задержки в шагах двигателя (она дублирует временную шкалу и нам не понадобится для анализа, поэтому мы будем её игнорировать).
    \item Прибор сохраняет эти файлы \textbf{без расширения}. Чтобы операционной системе и редакторам было удобнее с ними работать, мы добавили к ним стандартное расширение \texttt{.txt}.
    \item Имена файлов строго структурированы и несут в себе метаданные эксперимента, например: \texttt{10-0-1-bg\_1.txt}.
\end{itemize}

Давай разберем имя сигнального файла по частям, так как из него программа будет автоматически узнавать физические параметры замера:
\begin{enumerate}
    \item \texttt{10} — угол поворота первого поляризатора от горизонтальной оси (в градусах). В нашем эксперименте мы будем вращать именно его, исследуя угловую зависимость.
    \item \texttt{0} — угол поворота второго поляризатора (анализатора) в градусах. В данном цикле измерений он зафиксирован на нуле.
    \item \texttt{1} — номер повторного измерения (прохода, или реплики). Обычно делается несколько независимых последовательных проходов по времени задержки, чтобы потом математически снизить влияние случайных флуктуаций.
    \item \texttt{bg\_1} — имя (идентификатор) фонового референсного сигнала (сигнала без образца, проходящего через чистый воздух), на фоне которого измеряется образец в данном конкретном запуске.
\end{enumerate}

\subsection{Парсинг имен файлов}

Напишем функцию, которая будет автоматически разбирать (парсить) такое имя файла и возвращать все четыре параметра в виде удобного кортежа. Создай в папке \texttt{src} файл \texttt{data\_loader.py} и начни наполнять его кодом (полный листинг готового модуля на данном этапе также приведён в Приложении~\ref{app:data_loader}).

\begin{minipage}{\linewidth}
	\captionof{listing}{Парсинг имени экспериментального файла (data\_loader.py)}
	\begin{minted}{python}
from pathlib import Path
from utils import normalize_angle

def parse_filename(filename: str):
    """
    Извлекает параметры эксперимента из имени файла.
    Ожидаемый формат: 'angle1-angle2-repetition-bg_name.txt'
    """
    # Убираем расширение (.txt)
    name_without_ext = Path(filename).stem 
    
    # Разбиваем строку по дефису
    parts = name_without_ext.split('-')
    
    # Проверяем, что в имени именно 4 части
    if len(parts) != 4:
        raise ValueError(f"Неизвестный формат файла: {filename}")
        
    angle1 = float(parts[0])
    angle2 = float(parts[1])
    repetition = int(parts[2])
    bg_name = parts[3]
    
    # Нормализуем угол первого ротатора с помощью готовой утилиты из Раздела 2
    angle1 = normalize_angle(angle1)
        
    return angle1, angle2, repetition, bg_name
	\end{minted}
\end{minipage}

Если ты вызовешь эту функцию, передав ей тестовую строку с именем экспериментального файла:
\begin{minted}{python}
a1, a2, rep, bg = parse_filename("10-0-1-bg_1.txt")
print(a1, a2, rep, bg)
# Выведет: 10.0 0.0 1 bg_1
\end{minted}

\textbf{Давай разберём логику работы этой функции:}
\begin{itemize}
    \item Метод \texttt{Path(filename).stem} из стандартной библиотеки \texttt{pathlib} красиво отсекает расширение файла, оставляя чистую строку метаданных.
    \item Метод \texttt{split('-')} разрезает строку на список подстрок везде, где встречается знак дефиса.
    \item Приведение типов \texttt{float()} и \texttt{int()} превращает текстовые цифры в полноценные числовые переменные, готовые к математическим вычислениям.
    \item Импортированная функция \texttt{normalize\_angle(angle1)} из нашего собственного модуля утилит \texttt{utils.py} приводит сырой угол лимба к симметричному диапазону $[-180^\circ, 180^\circ]$. Это наглядный пример повторного использования кода: написав один раз качественную утилиту в Разделе 2, мы избавляем себя от дублирования громоздких ветвлений и связанных с ними багов в других модулях программы.
\end{itemize}

\warning{Обязательно делай защитные проверки, такие как \texttt{if len(parts) != 4:}, при разборе экспериментальных файлов. Если в папку с данными случайно попадет скрытый файл операционной системы или файл логов, программа не упадет посреди долгого расчета с непонятной ошибкой, а аккуратно сообщит о некорректном формате.}

\subsection{Загрузка массивов чисел с помощью NumPy}

Теперь, когда метаданные файла успешно прочитаны, нам нужно загрузить саму таблицу с числами временного сигнала. В научной среде общепринятым стандартом работы с массивами является библиотека \textbf{NumPy}. Она написана на языке C, благодаря чему обрабатывает огромные числовые колонки в сотни раз быстрее стандартных списков Python, и позволяет прочитать весь файл буквально одной строчкой.

Добавим функцию загрузки в наш файл \texttt{data\_loader.py}:

\begin{minipage}{\linewidth}
	\captionof{listing}{Загрузка экспериментальных данных (data\_loader.py)}
	\begin{minted}{python}
import numpy as np

def load_tds_data(filepath):
    """
    Загружает временной ТГц сигнал из текстового файла.
    Возвращает два одномерных массива NumPy: время (t) и амплитуду (E).
    """
    # unpack=True распаковывает колонки таблицы сразу в разные переменные.
    # usecols=(0, 1) указывает NumPy читать только первые две колонки,
    # полностью игнорируя третью колонку с позицией линии задержки.
    t, E = np.loadtxt(filepath, unpack=True, usecols=(0, 1))
    return t, E
	\end{minted}
\end{minipage}

Параметр \texttt{unpack=True} — это очень удобный прием в Python (так называемое распаковывание кортежа), который позволяет избежать промежуточного создания двумерной матрицы и сразу получить две независимые физические шкалы: шкалу времени $t$ и шкалу амплитуды $E$.

\subsection{Первичная визуализация сигналов: Знакомство с функцией plot}

В физике невозможно двигаться вслепую. После написания любой функции загрузки необходимо сразу же посмотреть на данные глазами. Для построения графиков в Python используется библиотека \textbf{Matplotlib}, а точнее её модуль \texttt{pyplot}, который предоставляет интуитивно понятный интерфейс для создания научных иллюстраций.

Давай проведем наш первый эксперимент и визуализируем сырой терагерцовый импульс. Открой файл \texttt{main.py} (пока он у нас выполняет роль проверочного скрипта) и напиши следующий код:

\begin{minipage}{\linewidth}
	\captionof{listing}{Визуализация ТГц-импульса во временной области (main.py)}
	\begin{minted}{python}
import matplotlib.pyplot as plt
import config
from data_loader import load_tds_data

# Строим путь к конкретному сигнальному файлу
test_file = config.DATA_DIR / "10-0-1-bg_1.txt"

# Загружаем числовые массивы
t, E = load_tds_data(test_file)

# Инициализируем графическое окно
plt.figure(figsize=(8, 4))

# Рисуем график временной зависимости
plt.plot(t, E, color='royalblue', linewidth=1.5, label='Сырой сигнал образца')

# Настраиваем оформление графика
plt.xlabel('Время задержки (пс)', fontsize=10)
plt.ylabel('Амплитуда поля E (у.е.)', fontsize=10)
plt.title('Терагерцовый импульс во временной области', fontsize=11, fontweight='bold')
plt.grid(True, linestyle='--', alpha=0.6)
plt.legend(loc='upper right')

# Отображаем окно на экране
plt.show()
	\end{minted}
\end{minipage}

Запусти файл \texttt{main.py}. На экране появится интерактивное окно с графиком ТГц-импульса (рис. \ref{fig:3_1}).

\smartfigure{fig3.1.png}{Типичный ТГц импульс во временной области после загрузки. Главный пик содержит основную энергию сигнала, а последующие осцилляции вызваны поглощением в водяном паре атмосферы.}{fig:3_1}

Как видишь, ТГц-сигнал во временной области представляет собой предельно короткий электромагнитный всплеск (субпикосекундной длительности), за которым следует плавный хвост из затухающих осцилляций.

\subsection{Физическое сравнение: Сигнал образца против фона (Референса)}

В терагерцовой спектроскопии мы никогда не анализируем сигнал образца изолированно. Измерения всегда проводятся парами:
\begin{enumerate}
    \item \textbf{Сигнал образца (sample):} ТГц-импульс проходит через исследуемую среду (в нашем случае — систему из двух поляризаторов и исследуемую решетку).
    \item \textbf{Фоновый сигнал (reference, или background):} ТГц-импульс проходит через то же самое оптическое плечо установки, но при полностью убранном образце (через чистый воздух).
\end{enumerate}

Давай напишем код, который загрузит оба файла (сигнал \texttt{50-0-1-bg\_4.txt} и его фоновый файл \texttt{bg\_4.txt}) и построит их на одном графике для наглядного физического сравнения.

\begin{minipage}{\linewidth}
	\captionof{listing}{Сопоставление сигнала и фона (main.py)}
	\begin{minted}{python}
# Загружаем сигнал образца
t_sig, E_sig = load_tds_data(config.DATA_DIR / "50-0-1-bg_4.txt")

# Загружаем соответствующий фоновый файл
t_bg, E_bg = load_tds_data(config.DATA_DIR / "bg_4.txt")

plt.figure(figsize=(9, 4.5))

# Строим две кривые на одной координатной сетке
plt.plot(t_bg, E_bg, color='crimson', linestyle='--', label='Опорный сигнал (Фон)', alpha=0.8)
plt.plot(t_sig, E_sig, color='navy', label='Сигнал образца', alpha=0.9)

plt.xlabel('Время задержки (пс)')
plt.ylabel('Амплитуда поля E (у.е.)')
plt.title('Сравнение ТГц импульсов: воздух vs образец')
plt.grid(True, linestyle=':')
plt.legend()
plt.show()
	\end{minted}
\end{minipage}

\smartfigure{fig3.2.png}{Сравнение временного ТГц-сигнала, прошедшего через образец, и опорного фонового сигнала (background). Наглядно видно затухание амплитуды из-за частичного поглощения и отражения.}{fig:3_2}

\note{\textbf{Физический анализ графиков:}\\
Внимательно посмотри на полученный график (рис. \ref{fig:3_2}). Мы можем сделать сразу важнейший физический вывод без всяких формул БПФ:\\
\textbf{Затухание (Амплитудный сдвиг):} Пиковое значение напряженности поля сигнала образца заметно меньше, чем у фона. Это говорит о том, что среда частично поглощает и отражает терагерцовое излучение.}

\subsection{Итоги раздела}
Мы проделали огромную инженерную и физическую работу:
\begin{enumerate}
    \item Разобрались в структуре экспериментальных файлов ТГц-установки и написали надежный парсер имен файлов.
    \item Научились использовать NumPy для быстрой поколоночной загрузки таблиц данных из текстовых файлов.
    \item Освоили библиотеку Matplotlib и построили наши первые графики терагерцовых импульсов, наглядно увидев физические эффекты затухания амплитуды при прохождении излучения через образец.
\end{enumerate}

Теперь мы научились загружать и визуализировать единичные файлы сигналов и фонов. Мы готовы перейти к самому сердцу цифровой обработки сигналов — расчету спектров пропускания с помощью быстрого преобразования Фурье. Этим мы займемся в следующем разделе.

